\documentclass[10pt,a4paper]{amsart}
\usepackage[margin=19mm]{geometry}
\usepackage{amsmath,amssymb,mathtools}
\usepackage[scr=rsfs,cal=euler]{mathalfa}
\usepackage{xfrac}
\usepackage[unicode,hidelinks,bookmarks=false]{hyperref}
\newcommand{\ie}{i.e.\ }
\newcommand{\cf}{cf.\ }
\newcommand{\resp}{respectively\ }
\newcommand{\Subsection}{\subsection}
\providecommand{\nobreakdash}{\nobreak\hbox{-}}
\setlength{\emergencystretch}{2em}
\multlinegap0pt
\newcommand{\EE}{\mathbb E}
\theoremstyle{plain}
\newtheorem{theorem}{Theorem}[section]
\newtheorem{proposition}[theorem]{Proposition}
\newtheorem{lemma}[theorem]{Lemma}
\newtheorem{corollary}[theorem]{Corollary}
\newtheorem{conjecture}[theorem]{Conjecture}
\theoremstyle{remark}
\newtheorem{remark}[theorem]{Remark}
\title[Rank-two Nahm sums, modularity obstructions, and the McKay--Thompson series $T_{31A}$]{Rank-two Nahm sums, modularity obstructions, and the McKay--Thompson series $T_{31A}$}
\author{Cetin Hakimoglu-Brown}
\date{Source: arXiv:2609.28509v1 (CC BY 4.0)}
\begin{document}
\maketitle

\noindent\emph{Source and licence.} The delivered work is the article shown in the title above, version arXiv:2609.28509v1 (CC BY 4.0), distributed under the Creative Commons Attribution 4.0 International licence, as recorded in the arXiv licence metadata for this exact version and shown on the abstract page saved with this item. The full licence notice, which carries the licence URL and the address of that abstract page, is the bundled file \texttt{licenses/arxiv-2609.28509-CC-BY-4.0.txt}.
\par\smallskip

\noindent\textit{Editorial scope.} Counted unit: the article's Lemma (source label lem:second) (source0.tex lines 834--840, source label \texttt{lem:second}): ``At the pair , 2 33813353 10556001 , 2- 12 1 2 - 1653155 3518667 .''. It is delivered together with the article's complete original proof (source0.tex lines 841--900, one proof environment adjacent to the statement, 60 lines), copied line for line from the article's own text. Required context retained uncounted: eq:local-log (lines 662--667); eq:reduction (lines 700--702); eq:wexplicit (lines 704--714); eq:gh (lines 720--727); eq:unique (lines 768--771); eq:L34 (lines 828--832). Editorial formatting only: an AMS \texttt{amsart} preamble that restates the article's own macro definitions and its own theorem declarations, and the theorem environments the retained lines use; the article's own class file is neither shipped nor input; the retained environments are renumbered sequentially in order of appearance, whereas the article prints them with its own numbering; the article's equation numbering is not reproduced and every cross-reference inside the retained text resolves to the wrapper's numbering. No citation occurs inside any retained line, so the wrapper carries no bibliography; All retained bodies are exact copies of the bundled \texttt{sources/211595/source0.tex}, so no omission receipt applies. No asset-inclusion command occurs inside any retained span and the package ships no image asset.


\section*{Retained context: eq:local-log (source0.tex lines 662--667)}

\begin{equation}\label{eq:local-log}
 \varepsilon^{-1}\Psi(t)+\log h(t)+\varepsilon c_0(t)
   +O(\varepsilon^3),\qquad
 c_0(t)=-\sum_{i=1}^2d_i
       \left(\frac1{24}+\frac{e^{-t_i}}{12(1-e^{-t_i})}\right).
\end{equation}


\section*{Retained context: eq:reduction (source0.tex lines 700--702)}

\begin{equation}\label{eq:reduction}
 \rho^3=5\rho^2-2\rho-1.
\end{equation}


\section*{Retained context: eq:wexplicit (source0.tex lines 704--714)}

\begin{equation}\label{eq:wexplicit}
 w_1=1+4\rho-\rho^2,\qquad
 w_2=774+2546\rho-603\rho^2.
\end{equation}
\begin{equation}\label{eq:Cexplicit}
 \Sigma=\frac1{361}
 \begin{pmatrix}
 -566+949\rho-222\rho^2&3596-5611\rho+1271\rho^2\\
 3596-5611\rho+1271\rho^2&-20832+32302\rho-7099\rho^2
 \end{pmatrix}.
\end{equation}


\section*{Retained context: eq:gh (source0.tex lines 720--727)}

\begin{align}
 g_i&=-\frac{b_i}{d_i}-\frac{w_i}{2},&
 h_i&=\frac{w_i(1+w_i)}2,\label{eq:gh}\\
 \psi_{3,i}&=\frac{w_i(1+w_i)}{d_i},&
 \psi_{4,i}&=-\frac{w_i(1+w_i)(1+2w_i)}{d_i},\label{eq:TV}\\
 c_0&=-\sum_{i=1}^2d_i\left(\frac1{24}+\frac{w_i}{12}\right).
 \label{eq:c0}
\end{align}


\section*{Retained context: eq:unique (source0.tex lines 768--771)}

\begin{equation}\label{eq:unique}
 (a,b)=\left(\frac{19}{3},\frac{124}{3}\right),
 \qquad \beta_1=\frac{8806}{3249}.
\end{equation}


\section*{Retained context: eq:L34 (source0.tex lines 828--832)}

\begin{align}
 L_3&=\sum_i(c_i'\xi_i+j_i\xi_i^3/6+\psi_{5,i}\xi_i^5/120),\notag\\
 L_4&=\sum_i(c_i''\xi_i^2/2+k_i\xi_i^4/24+\psi_{6,i}\xi_i^6/720).
 \label{eq:L34}
\end{align}


\section{The counted unit: Lemma (source label lem:second) and its complete original proof}

\begin{lemma}\label{lem:second}
At the pair \eqref{eq:unique},
\begin{equation}\label{eq:beta2value}
 \beta_2=\frac{33813353}{10556001},\qquad
 \beta_2-\frac12\beta_1^2=-\frac{1653155}{3518667}.
\end{equation}
\end{lemma}

\begin{proof}
The expansion of the exponential of \eqref{eq:local-log} gives
\begin{equation}\label{eq:beta2gauss}
 \beta_2=
 \EE\left(L_4+L_1L_3+\frac12L_2^2
                 +\frac12L_1^2L_2+\frac1{24}L_1^4\right).
\end{equation}
Indeed, after removal of the leading Gaussian, the exponent is
$\sqrt\varepsilon L_1+\varepsilon L_2+
\varepsilon^{3/2}L_3+\varepsilon^2L_4+\cdots$.
Odd Gaussian moments vanish. The five terms in
\eqref{eq:beta2gauss} are exactly the remaining contributions of
order $\varepsilon^2$.

Here is a finite algebraic specification of the evaluation.
At \eqref{eq:unique} the two linear coefficients are
\[
 g_1=\frac{\rho^2}{2}-2\rho-\frac{41}{6},\qquad
 g_2=\frac{603\rho^2}{2}-1273\rho-\frac{1165}{3}.
\]
Use these values, \eqref{eq:wexplicit}--\eqref{eq:Cexplicit}, and
\eqref{eq:gh}--\eqref{eq:L34}.
For a polynomial $p(u,v)$ let
\[
 \mathfrak E_\Sigma(p)=
 \left.
 \sum_{\ell=0}^{6}\frac1{2^\ell\ell!}
 \left(\Sigma_{11}\partial_u^2+
       2\Sigma_{12}\partial_u\partial_v+
       \Sigma_{22}\partial_v^2\right)^\ell p(u,v)
 \right|_{u=v=0}.
\]
This is Gaussian expectation for polynomials of degree at most twelve,
which includes every term of \eqref{eq:beta2gauss}.
Consequently the exact arithmetic identity used here is
\begin{equation}\label{eq:certificate}
 \mathfrak E_\Sigma\left(
 L_4+L_1L_3+\tfrac12L_2^2+
 \tfrac12L_1^2L_2+\tfrac1{24}L_1^4
 \right)
 =\frac{33813353}{10556001}\qquad\text{in }K.
\end{equation}
Here $\xi_1=u$ and $\xi_2=v$. All entries are given explicitly above; only polynomial
multiplication, differentiation, and the reduction
\eqref{eq:reduction} are involved. An equivalent way to evaluate the
left side is to use $M_{00}=1$, zero moments of odd total degree,
and the recurrence
\begin{align}
 M_{r,s}&=(r-1)\Sigma_{11}M_{r-2,s}
              +s\Sigma_{12}M_{r-1,s-1}\quad(r>0),\notag\\
 M_{0,s}&=(s-1)\Sigma_{22}M_{0,s-2}\quad(s>0),
 \label{eq:momentrecurrence}
\end{align}
where negative indices give zero.
The reduced coefficients of $1,\rho,\rho^2$ in the left side of
\eqref{eq:certificate} are respectively
$33813353/10556001$, $0$, and $0$.
Finally, $\beta_1=8806/3249$ and $3249^2=10556001$ give
the second identity in \eqref{eq:beta2value}.
\end{proof}

\end{document}
